\chapter{BV approximation and compactness}\label{sec:bv-approximation}
We now develop the analytic compactness behind finite-perimeter minimization. Here $U\subset\mathbb R^n$ is open, and $Du$ denotes the distributional vector measure of Definition \ref{def:bv}. The basic reference is \cite{Sim18}, Chapter 2, Section 2, pp.~52--56. The proofs below spell out the mollification and compactness steps needed in the direct method.

\begin{proposition}[Lower semicontinuity of total variation]\label{prop:bv-lsc}
If $u_j\to u$ in $L^1_{\mathrm{loc}}(U)$, then for every open $W\subset U$,
\[
 |Du|(W)\leq\liminf_j|Du_j|(W),
\]
using the variational definition even when its value is infinite. In particular, uniform local variation bounds force $u\in BV_{\mathrm{loc}}(U)$.
\end{proposition}
\begin{proof}
For $X\in C_c^1(W;\mathbb R^n)$ with $|X|\leq1$, local $L^1$ convergence gives $\int u\operatorname{div}X=\lim_j\int u_j\operatorname{div}X\leq\liminf_j|Du_j|(W)$. Take the supremum over $X$. This is the general-function form of Proposition \ref{prop:perimeter-lsc}, and the lower-semicontinuity step in \cite{Sim18}, Chapter 2, Theorem 2.6.
\end{proof}

\begin{proposition}[Multiplication by a smooth cutoff]\label{prop:bv-cutoff-product}
If $u\in BV_{\mathrm{loc}}(U)$ and $\eta\in C_c^1(U)$, then the function $\eta u$, extended by zero outside $U$, belongs to $BV(\mathbb R^n)$ and
\[
 D(\eta u)=\eta Du+uD\eta\,\mathcal L^n,
 \qquad |D(\eta u)|(\mathbb R^n)
 \leq\int|\eta|\,d|Du|+\int|u||D\eta|\,dx.
\]
The compact support of the cutoff prevents a boundary measure from appearing at $\partial U$.
\end{proposition}
\begin{proof}
For a compactly supported test field $X$, use $\eta\operatorname{div}X=\operatorname{div}(\eta X)-D\eta\cdot X$. The defining identity for $Du$ then gives the asserted vector measure. Its total variation is bounded by the sum of the total variations of the two terms. Both are finite on the compact support, and $\eta u$ is integrable. Compare the product estimate in \cite{Sim18}, Chapter 2, the proof of Lemma 2.8.
\end{proof}

\begin{lemma}[Mollification and the variation bound]\label{lem:bv-mollification}
Let $u\in BV(\mathbb R^n)$ and let $\rho\in C_c^\infty(B_1)$ be nonnegative with integral one. Put $\rho_\varepsilon(x)=\varepsilon^{-n}\rho(x/\varepsilon)$ and $u_\varepsilon=\rho_\varepsilon*u$. Then $u_\varepsilon\to u$ in $L^1(\mathbb R^n)$ and
\[
 Du_\varepsilon(x)=\int\rho_\varepsilon(x-y)\,dDu(y),
 \qquad |Du_\varepsilon(x)|\leq\int\rho_\varepsilon(x-y)\,d|Du|(y).
\]
In particular, $\int|Du_\varepsilon|\leq|Du|(\mathbb R^n)$, and for a Borel set $A$, $\int_A|Du_\varepsilon|\leq|Du|(A_\varepsilon)$, where $A_\varepsilon=\{y:\operatorname{dist}(y,A)<\varepsilon\}$. See \cite{Sim18}, Chapter 2, Lemma 2.5.
\end{lemma}
\begin{proof}
Convolution is an approximate identity on $L^1$, giving the first convergence. Distributional differentiation commutes with convolution: test against a compactly supported smooth function, use Fubini, and then apply the defining integration-by-parts identity for $Du$. The triangle inequality for integration against a vector measure gives the pointwise bound. Integrate it and use Fubini again. The kernel has total integral one, and it vanishes when $|x-y|\geq\varepsilon$, which proves the two integral bounds.
\end{proof}

\begin{definition}[Strict convergence in BV]\label{def:bv-strict-convergence}
A sequence $u_j\in BV(U)$ converges strictly to $u\in BV(U)$ if
\[
 \|u_j-u\|_{L^1(U)}\to0,
 \qquad |Du_j|(U)\to|Du|(U).
\]
This does not mean $|D(u_j-u)|(U)\to0$. In particular, convergence of the total variations is different from convergence in the full BV norm. The terminology describes the approximation property in \cite{Sim18}, Chapter 2, Lemma 2.5; the local Radon-measure version there requires localization.
\end{definition}

\begin{corollary}[Global mollification is strict approximation]\label{cor:bv-strict-mollification}
The mollifications in Lemma \ref{lem:bv-mollification} converge strictly to $u$ in $BV(\mathbb R^n)$.
\end{corollary}
\begin{proof}
The lemma gives the $L^1$ convergence and the upper bound $\limsup\int|Du_\varepsilon|\leq|Du|(\mathbb R^n)$. Proposition \ref{prop:bv-lsc} gives the reverse lower bound. Both together yield convergence of the total variations. On a general domain, arbitrary zero extension can create a boundary jump; the whole-space hypothesis is part of this statement.
\end{proof}

\begin{proposition}[Translation estimate for BV functions]\label{prop:bv-translations}
For $u\in BV(\mathbb R^n)$ and $h\in\mathbb R^n$,
\[
 \int_{\mathbb R^n}|u(x+h)-u(x)|\,dx
 \leq |h|\,|Du|(\mathbb R^n).
\]
Consequently, with the mollifier from Lemma \ref{lem:bv-mollification},
\[
 \|u_\varepsilon-u\|_{L^1}\leq
 \varepsilon\left(\int|z|\rho(z)\,dz\right)|Du|(\mathbb R^n).
\]
\end{proposition}
\begin{proof}
For the smooth function $u_\delta$, the fundamental theorem of calculus along $x+th$, followed by Fubini, gives
\[
 \int|u_\delta(x+h)-u_\delta(x)|\,dx
 \leq|h|\int_0^1\int|Du_\delta(x+th)|\,dx\,dt
 \leq|h||Du|(\mathbb R^n).
\]
Pass to the limit $\delta\downarrow0$ in $L^1$. Average the translation estimate against $\rho_\varepsilon(h)$ to get the second inequality. This supplies a quantitative form of the smoothing step in \cite{Sim18}, Chapter 2, Theorem 2.6.
\end{proof}

\begin{theorem}[Local compactness of BV functions]\label{thm:bv-local-compactness}
Suppose $u_j\in BV_{\mathrm{loc}}(U)$ and for every relatively compact open $W\subset U$,
\[
 \sup_j\left(\|u_j\|_{L^1(W)}+|Du_j|(W)\right)<\infty.
\]
Then a subsequence converges in $L^1_{\mathrm{loc}}(U)$ to a function $u\in BV_{\mathrm{loc}}(U)$. The lower-semicontinuity conclusion of Proposition \ref{prop:bv-lsc} holds on every open subset. This is \cite{Sim18}, Chapter 2, Theorem 2.6, pp.~54--55.
\end{theorem}
\begin{proof}
First suppose the functions are globally BV, supported in one compact set $K$, and have uniformly bounded $L^1$ norm and total variation. For fixed $\varepsilon>0$, their convolutions with $\rho_\varepsilon$ have common compact support in $K+\overline{B_\varepsilon}$ and uniform bounds for their values and first derivatives, obtained from the $L^1$ bound and the suprema of $\rho_\varepsilon$ and $D\rho_\varepsilon$. Arzel\`a--Ascoli therefore makes these convolutions precompact in $L^1$.

Proposition \ref{prop:bv-translations} bounds the $L^1$ distance from each function to its convolution by $C\varepsilon$, uniformly in $j$. For each prescribed tolerance, the convolutions have a finite $L^1$ net, which becomes a finite net for the original functions with twice the error after choosing $\varepsilon$ small. Hence the original sequence is precompact in $L^1$.

For the local statement, choose an exhaustion $W_k\Subset U$ and cutoffs $\eta_k\in C_c^1(U)$ equal to one on $W_k$. Proposition \ref{prop:bv-cutoff-product} gives uniform whole-space BV bounds for $\eta_ku_j$, with a common compact support for each $k$. Extract subsequences successively and then take a diagonal subsequence. The limits agree on overlaps because they are limits in $L^1$, and therefore define $u\in L^1_{\mathrm{loc}}(U)$. Proposition \ref{prop:bv-lsc} gives local finite variation and the asserted bounds.
\end{proof}

\begin{corollary}[Compactness of finite-perimeter sets]\label{cor:perimeter-compactness}
Let $E_j$ be measurable sets with uniformly bounded perimeter on every relatively compact open subset of $U$. Then a subsequence satisfies $\chi_{E_j}\to\chi_E$ in $L^1_{\mathrm{loc}}(U)$ for a set $E$ of locally finite perimeter. If all $E_j$ lie almost everywhere in one compact set $K\subset\mathbb R^n$ and have uniformly bounded total perimeter, the convergence is in $L^1(\mathbb R^n)$, the limit also lies in $K$ almost everywhere, and the volumes converge.
\end{corollary}
\begin{proof}
On a relatively compact $W$, the characteristic functions have $L^1$ norm at most $\mathcal L^n(W)$, so Theorem \ref{thm:bv-local-compactness} applies. A further subsequence converges almost everywhere on an exhaustion; since every value is zero or one, its limit is a characteristic function. Proposition \ref{prop:perimeter-lsc} gives the perimeter bound. With common compact support, use the whole-space compactness argument in the theorem. Global $L^1$ convergence implies convergence of the integrals and hence of the volumes.
\end{proof}

\begin{proposition}[Weak convergence of the derivative measures]\label{prop:bv-derivative-convergence}
If $u_j\to u$ in $L^1_{\mathrm{loc}}(U)$ and their variations are uniformly bounded on relatively compact open sets, then $Du_j\rightharpoonup Du$ as vector Radon measures locally in $U$. Convergence of the variation measures $|Du_j|$ to $|Du|$ is not asserted.
\end{proposition}
\begin{proof}
Testing with a smooth compactly supported vector field $X$ gives $\int X\cdot dDu_j=-\int u_j\operatorname{div}X\to-\int u\operatorname{div}X$. A continuous compactly supported field can be approximated uniformly by smooth ones supported in a fixed relatively compact neighborhood. The variation bound controls the error uniformly in $j$, proving weak convergence against every continuous test field. This is the distributional-measure argument of \cite{Sim18}, Chapter 2, Section 2.
\end{proof}

\begin{example}[Variation control without function control is insufficient]\label{ex:bv-constants}
On a bounded nonempty open set, the constants $u_j=j$ have $|Du_j|=0$, but no subsequence converges in $L^1_{\mathrm{loc}}$. The local $L^1$ hypothesis in Theorem \ref{thm:bv-local-compactness} cannot be omitted. For characteristic functions this hypothesis is automatic locally, which explains the simpler assumptions in Corollary \ref{cor:perimeter-compactness}.
\end{example}

\begin{example}[Strict convergence need not be BV-norm convergence]\label{ex:bv-strict-not-norm}
In one dimension let $u=\chi_{(0,1)}$ and let $u_\varepsilon$ be its nonnegative smooth mollifications. Then $u_\varepsilon\to u$ strictly by Corollary \ref{cor:bv-strict-mollification}, and $|Du|(\mathbb R)=2$. But $Du=\delta_0-\delta_1$ is singular with respect to Lebesgue measure, whereas $Du_\varepsilon$ is absolutely continuous. Hence these derivative measures are mutually singular and
\[
 |D(u_\varepsilon-u)|(\mathbb R)=\int|Du_\varepsilon|\,dx+2\longrightarrow4.
\]
Thus the stronger BV norm does not tend to zero. This example also distinguishes weak convergence of vector measures from convergence of their total variation norms.
\end{example}

\chapter{Perimeter blow-ups and the direct method}\label{sec:perimeter-direct-method}
Compactness and lower semicontinuity establish existence before smoothness is known. The reduced-boundary blow-up identifies the local geometric object, but by itself does not make every topological boundary point regular. In this section $n\geq2$.

\begin{definition}[Measure-theoretic interior and exterior]\label{def:measure-theoretic-boundary}
For a measurable set $E\subset\mathbb R^n$, define
\[
 E^{(1)}=\left\{x:\lim_{r\downarrow0}\frac{|E\cap B_r(x)|}{|B_r(x)|}=1\right\},
 \qquad
 E^{(0)}=\left\{x:\lim_{r\downarrow0}\frac{|E\cap B_r(x)|}{|B_r(x)|}=0\right\}.
\]
The essential boundary is $\partial^eE=\mathbb R^n\setminus(E^{(0)}\cup E^{(1)})$. These definitions are unchanged by a Lebesgue-null modification of $E$. They differ from the topological interior, exterior, and boundary, which depend on the chosen representative. Compare the measure-density framework of \cite{Sim18}, Chapters 1 and 3.
\end{definition}

\begin{theorem}[Half-space blow-up at a reduced-boundary point]\label{thm:perimeter-halfspace-blowup}
Let $E$ have locally finite perimeter in $U$, and let $x\in\partial^*E\cap U$. With outward normal $\nu_E(x)$, set $E_{x,r}=(E-x)/r$. As $r\downarrow0$,
\[
 \chi_{E_{x,r}}\longrightarrow\chi_{\{y:y\cdot\nu_E(x)<0\}}
 \quad\text{in }L^1_{\mathrm{loc}}(\mathbb R^n),
\]
and the perimeter measures converge locally to $\mathcal H^{n-1}\llcorner\nu_E(x)^\perp$. In particular,
\[
 \lim_{r\downarrow0}\frac{|E\cap B_r(x)|}{\omega_nr^n}=\frac12,
 \qquad
 \lim_{r\downarrow0}\frac{P(E;B_r(x))}{\omega_{n-1}r^{n-1}}=1.
\]
This is the blow-up conclusion in the proof of \cite{Sim18}, Chapter 3, Theorem 4.4, especially equations (9)--(13), pp.~98--100, with the project's outward-normal convention. The substantial blow-up and density argument remains external. The convergence holds as $r\downarrow0$, not merely for one selected subsequence.
\end{theorem}

\begin{example}[A null modification can change the whole topological boundary]\label{ex:null-modification-boundary}
The set $E=\mathbb Q^n$ is null and dense, and its complement is also dense. Thus its topological boundary is all of $\mathbb R^n$, while $\chi_E=0$ almost everywhere, $P(E;U)=0$, $E^{(0)}=\mathbb R^n$, and its reduced and essential boundaries are empty. Perimeter measures are therefore attached to measurable sets modulo null sets, not to arbitrary topological representatives.
\end{example}

\begin{proposition}[Existence with a fixed volume in a compact container]\label{prop:perimeter-volume-existence}
Let $K\subset\mathbb R^n$ be compact and $v\geq0$. Suppose there is at least one measurable $E\subset K$ almost everywhere with $|E|=v$ and finite total perimeter. Then the infimum of $P(E;\mathbb R^n)$ over this class is attained. No regularity of $K$ is needed for this existence statement, and no smoothness of the minimizing boundary is concluded.
\end{proposition}
\begin{proof}
Choose a minimizing sequence with bounded total perimeter. All its characteristic functions are supported in $K$ and have fixed $L^1$ norm $v$. Corollary \ref{cor:perimeter-compactness} gives a globally $L^1$-convergent subsequence whose limit remains supported in $K$ and has volume $v$. Proposition \ref{prop:perimeter-lsc} shows that this limit attains the infimum. This is a direct application of the BV compactness theorem in \cite{Sim18}, Chapter 2.
\end{proof}

\begin{proposition}[A local direct method with a bulk term]\label{prop:perimeter-bulk-existence}
Let $U\subset\mathbb R^n$ be bounded and open, let $K\Subset U$ be compact, let $E_0$ have finite perimeter in $U$, and let $g\in L^1(U)$. Then
\[
 \mathcal J_g(E)=P(E;U)-\int_{E\cap U}g(x)\,dx
\]
attains its infimum among measurable sets $E$ with $\chi_E=\chi_{E_0}$ almost everywhere on $U\setminus K$. The compactly contained set of allowed changes is part of the boundary-data formulation. This extends the constant bulk term in Definition \ref{def:cmc-functional}.
\end{proposition}
\begin{proof}
The functional is bounded below by $-\|g\|_{L^1(U)}$, and $E_0$ is an admissible finite-energy competitor. A minimizing sequence therefore has uniformly bounded perimeter. Corollary \ref{cor:perimeter-compactness} gives a subsequence converging locally in $L^1$ and, after a further extraction, almost everywhere. The limit agrees with $E_0$ on $U\setminus K$. Dominated convergence, using $|g\chi_E|\leq|g|$, shows convergence of the bulk integrals. Perimeter lower semicontinuity then proves attainment. The argument uses no trace-continuity assertion at $\partial U$, because every competitor is fixed outside $K$.
\end{proof}

\begin{proposition}[The Euler equation on a smooth minimizing interface]\label{prop:perimeter-bulk-euler}
Suppose $g\in C^1(U)$ and $E$ minimizes the functional in Proposition \ref{prop:perimeter-bulk-existence}. On any open region $W\subset\operatorname{int}K$ where $E$ is represented by a domain with smooth boundary, that boundary satisfies
\[
 H=g
\]
with respect to the outward unit normal and the trace convention $H=\operatorname{div}_{\partial E}\nu$. This is an interior conclusion; it makes no assertion about contact with the container or about singular points.
\end{proposition}
\begin{proof}
Use the flow of a smooth vector field compactly supported in $W$ to vary $E$. Both positive and negative variation times are admissible. If $f=X\cdot\nu$ is the normal speed, the first variation of perimeter is $\int_{\partial E}Hf\,dA$, while the transport formula for the bulk integral gives $\int_{\partial E}gf\,dA$. Minimality makes their difference zero. Every compactly supported smooth normal speed on the interface can be extended to an ambient field, so $H-g=0$ there. Compare Proposition \ref{prop:cmc-first-variation} for constant $g$.
\end{proof}
